\begin{question}{Efficiently Simulating Bounded Stacks with RNNs}

\begin{subquestion}
    \begin{table}[H]
    \centering
        \begin{tabular}{cccc}
            \toprule
            & $\stackseq_1 = \openBr{2}$ &
            $\stackseq_2 = \openBr{1} \openBr{1} \openBr{1}$ &
            $\stackseq_3 = \openBr{2} \openBr{1} \openBr{2}$ \\
            \midrule
            $\closeBr{2}$ & & &\\
            % \midrule
            $\closeBr{1}$ & & & \\
            % \midrule
            $\openBr{2}$ & & & \\
            \bottomrule
        \end{tabular}
        \caption{The first row of the table represents different stack configurations, while the first column lists the newly read symbols.}
        \label{tab:stackconfigs}
    \end{table}
\end{subquestion}

\begin{subquestion}

\end{subquestion}

\begin{subquestion}
    \begin{table}[H]
    \centering
        \begin{tabular}{cccc}
            \toprule
            & $\stackseq_1 = \openBr{2}$ &
            $\stackseq_2 = \openBr{1} \openBr{1} \openBr{1}$ &
            $\stackseq_3 = \openBr{2} \openBr{1} \openBr{2}$ \\
            \midrule
            $\closeBr{2}$ & & &\\
            % \midrule
            $\closeBr{1}$ & & & \\
            % \midrule
            $\openBr{2}$ & & & \\
            \bottomrule
        \end{tabular}
        \caption{The first row of the table represents different stack configurations, while the first column lists the newly read symbols.}
        \label{tab:stackconfigs}
    \end{table}
\end{subquestion}

\begin{subquestion}

\end{subquestion}

\begin{subquestion}

\end{subquestion}

\begin{subquestion}

\end{subquestion}

\begin{subquestion}

\end{subquestion}

\begin{subquestion}

\end{subquestion}

\begin{subquestion}

\end{subquestion}

\end{question}